\bookchapter{Capítulo 55: Projeto 2 — App de consulta de clima (versão conectada)}

\chapterintro{O controle de gastos é 'treinar o fundamento'; o App de clima é 'ver o mundo' — ele precisa lidar com a internet. Neste projeto você vai aprender: pedir uma API gratuita, fazer parsing de JSON, tratar os três estados de interface 'carregando/sucesso/erro' e como esconder bem a chave da API. Ao terminá-lo, você terá a capacidade de fazer qualquer 'App conectado' — notícias, câmbio, tradução; a receita é a mesma.}

\booksection{Pedindo uma API de clima gratuita}

Para estudo individual, recomendo o OpenWeatherMap ou, na China, a QWeather — ambas têm cota gratuita e funcionam assim que você se cadastra. Usando o OpenWeatherMap como exemplo:

\begin{enumerate}
  \item Crie a conta na página API keys, gere uma chave (apikey);
  \item O plano gratuito permite chamar o endpoint Current Weather (clima atual): api.openweathermap.org/data/2.5/weather? q=beijing\&appid=SUA\_CHAVE\&units=metric;
  \item Teste primeiro no navegador para ver como é o JSON que volta — esse 'preparo' é obrigatório antes de integrar qualquer API.
\end{enumerate}

\textbf{Alerta de armadilha: a chave jamais pode ficar escrita no código}\par

A API Key é como a 'senha da sua carteira': colocada no código e enviada ao GitHub, em poucas horas os crawlers já a terão varrido para uso fraudulento. O jeito certo: coloque-a em local.properties (arquivo já ignorado pelo gitignore) e leia-a no código via buildConfig do Gradle. Pesquise 'melhores práticas para ocultar a API Key no Android' e siga as instruções.

\booksection{Camada de rede: Retrofit + parsing de JSON}

Primeiro defina o 'modelo de dados' conforme o JSON que a API devolve, ignorando todos os campos que não casam:

\begin{lstlisting}[language=Kotlin]
// JSON: {"weather":[{"description":"Ensolarado"}],"main":{"temp":26.5,"humidity":40},"name":"Pequim"}
data class WeatherResponse(
    val weather: List<Weather>,
    val main: Main,
    val name: String
)
data class Weather(val description: String)
data class Main(val temp: Double, val humidity: Int)

interface WeatherApi {
    @GET("data/2.5/weather")
    suspend fun getWeather(
        @Query("q") city: String,
        @Query("appid") key: String,
        @Query("units") metric: String = "metric",
        @Query("lang") lang: String = "pt"
    ): WeatherResponse
}
\end{lstlisting}

O Gson enche esses data class com o JSON automaticamente — com os nomes dos campos iguais, a mágica acontece sozinha.

\booksection{Os três estados de interface: carregando / sucesso / erro}

A alma de um App conectado está no 'gerenciamento de estados'. Rede lenta, internet caindo, nome de cidade errado... em cada caso, a interface precisa dar uma resposta ao usuário. Represente isso com elegância usando um sealed class (classe selada):

\begin{lstlisting}[language=Kotlin]
sealed class UiState {
    object Loading : UiState()                    // carregando
    data class Success(val data: WeatherResponse) : UiState()  // sucesso
    data class Error(val message: String) : UiState()          // errore
}

// No ViewModel:
fun search(city: String) {
    viewModelScope.launch {
        _state.value = UiState.Loading
        _state.value = try {
            UiState.Success(api.getWeather(city))
        } catch (e: Exception) {
            UiState.Error("Não encontrei essa cidade — tentar outro nome?")
        }
    }
}

// Na tela:
when (val s = state) {
    is UiState.Loading -> CircularProgressIndicator()   // carregando
    is UiState.Success -> WeatherContent(s.data)        // exibe
    is UiState.Error -> Text(s.message)                 // erro
}
\end{lstlisting}

Esse formato em três partes Loading/Success/Error é o 'modelo curinga' das interfaces conectadas: memorize-o e, depois, copie a fórmula em todos os Apps conectados que fizer.

\booksection{Código de interface na prática: desenhando o clima}

Com os estados prontos, a interface fica assim (trechos principais):

\begin{lstlisting}[language=Kotlin]
@Composable
fun WeatherContent(data: WeatherResponse) {
    Column(
        Modifier.fillMaxSize().padding(24.dp),
        horizontalAlignment = Alignment.CenterHorizontally
    ) {
        Text(data.name, fontSize = 28.sp, fontWeight = FontWeight.Bold)
        Text(
            "${data.main.temp}°C",
            fontSize = 72.sp,
            fontWeight = FontWeight.Light
        )
        Text(
            data.weather.firstOrNull()?.description ?: "Desconhecido",
            fontSize = 20.sp,
            color = Color.Gray
        )
        Spacer(Modifier.height(16.dp))
        Row {
            Text(" Umidade ${data.main.humidity}%")
            Spacer(Modifier.width(24.dp))
            Text(" Sensação ${data.main.temp}°C")
        }
    }
}
\end{lstlisting}

Alguns detalhes que valem a pena aprender: firstOrNull() pega o primeiro elemento com segurança (a lista pode estar vazia) e ?: “Desconhecido” oferece um texto de reserva — o princípio da null safety vale na interface também. A temperatura usa fonte grande de 72.sp para dar presença, e o peso 'leve' (Light) deixa os números elegantes; essa é a receita clássica dos Apps de clima.

\booksection{Detalhes de interface e interação}

\booksection{Detalhes de interface e interação}

\begin{itemize}
  \item Campo de busca: TextField no topo; digite a cidade e toque em buscar (configure a ação 'buscar' no teclado — dar Enter já pesquisa);
  \item Cartão do clima: temperatura em fonte grande (26,5°), descrição do tempo (céu limpo), umidade e nome da cidade;
  \item Cidades favoritas: salve no DataStore as cidades que você consulta sempre e consulte-as automaticamente ao reabrir — de quebra, uma revisão do capítulo 6;
  \item Puxar para atualizar: no Compose, use o PullToRefreshBox (oferecido pelo Material 3) e a sensação fica 'profissional' na hora.
\end{itemize}

\begin{tcbinfo}{Dica: imagens dão vida à interface}
Dê a cada tipo de clima um emoji ou ícone () mapeado com when(description) — três linhas de código e o charme da interface dobra. Também dá para usar uma biblioteca de ícones de clima de terceiros, mas na fase de estudo os emojis bastam.

\end{tcbinfo}

\begin{tcbpractice}{Mãos à obra: checklist de aceitação do App de clima}
(1) A busca de clima funciona normalmente; (2) sem internet, o App comunica o erro com elegância sem cair (teste com o modo avião ativado!); (3) digitar uma cidade inexistente gera um aviso amigável; (4) favoritar duas cidades e, ao reiniciar o App, elas continuam lá; (5) girar a tela não perde o estado. Só com as cinco aprovadas o projeto se forma.

\end{tcbpractice}

\begin{tcbsummary}{Resumo do capítulo}
\begin{itemize}
  \item Teste a API no navegador e entenda o JSON antes de escrever o código — 'se preparar antes' nunca dá prejuízo
  \item A API Key vai escondida em local.properties; jamais entra no repositório de código
  \item Interface Retrofit + data class: com os campos casados, o JSON é convertido automaticamente
  \item Os três estados Loading/Success/Error são o curinga das interfaces conectadas
  \item Sem internet, entrada vazia, cidade inexistente — tratar erros é a nota mínima de um App conectado
\end{itemize}

\end{tcbsummary}
